\section{Integral monodromy and the base orbifold}

\subsection{The lattice representation}
Let $V=\Z^4$ with ordered basis $(\gamma,u,w,\delta)$ and let
$\Lambda=V^*$ with dual basis
$(\wh\gamma,\wh u,\wh w,\wh\delta)$.  Evaluation on $\gamma$ is denoted
by the same letter:
\[
 \gamma(a\wh\gamma+b\wh u+c\wh w+d\wh\delta)=a.
\]
Set
\[
T_1=
\begin{pmatrix}
1&0&-6&2\\
0&-1&1&1\\
0&-1&0&1\\
0&0&0&1
\end{pmatrix},
\qquad
T_2=
\begin{pmatrix}
1&6&0&-3\\
0&0&-1&1\\
0&1&0&0\\
0&0&0&1
\end{pmatrix}.
\]
For $T\in\GL(V)$ write $A(T)=(T^{-1})^t\in\GL(\Lambda)$, and put
$A_j=A(T_j)$.

\begin{proposition}[Integral linear algebra]\label{prop:linear-algebra}
The following statements hold.
\begin{enumerate}[label=\textup{(\roman*)}]
\item $T_1$ and $T_2$ have exact orders $3$ and $4$, respectively, and
\[
T_0:=(T_1T_2)^{-1}
 =\begin{pmatrix}
1&0&0&1\\
0&1&-1&0\\
0&0&1&0\\
0&0&0&1
\end{pmatrix}=I+N,
\qquad N^2=0,
\]
with
\[
N\gamma=Nu=0,\qquad Nw=-u,\qquad N\delta=\gamma.
\]
\item In the dual basis,
\[
\begin{aligned}
A_1\wh\gamma&=\wh\gamma+6\wh u-6\wh w-2\wh\delta,
&A_1\wh u&=-\wh w+\wh\delta,
&A_1\wh w&=\wh u-\wh w,
&A_1\wh\delta&=\wh\delta,\\
A_2\wh\gamma&=\wh\gamma-6\wh w+3\wh\delta,
&A_2\wh u&=\wh w,
&A_2\wh w&=-\wh u+\wh\delta,
&A_2\wh\delta&=\wh\delta.
\end{aligned}
\]
Writing $M_0=A(T_0)$, one has
\[
M_0\wh\gamma=\wh\gamma-\wh\delta,
\quad M_0\wh u=\wh u+\wh w,
\quad M_0\wh w=\wh w,
\quad M_0\wh\delta=\wh\delta,
\]
and $A_1A_2M_0=I$.
\item Put
\[
 \Lambda_{\mathrm{tor}}=\langle\wh w,\wh\delta\rangle,
 \qquad \ol\Lambda=\Lambda/\Lambda_{\mathrm{tor}}.
\]
Then
\[
\ker(M_0-I)=\im(M_0-I)=\Lambda_{\mathrm{tor}},
\]
and the induced map
\[
 B_0:\ol\Lambda\longrightarrow\Lambda_{\mathrm{tor}},
 \qquad B_0\ol{\wh\gamma}=-\wh\delta,
 \quad B_0\ol{\wh u}=\wh w,
\]
is unimodular.
\item The vectors
\[
 \eps=\wh\gamma+2\wh u-4\wh w,
 \qquad
 \eps'=\wh\gamma+3\wh u-3\wh w
\]
satisfy
\[
 \Lambda^{A_1}=\langle\eps,\wh\delta\rangle,
 \qquad
 \Lambda^{A_2}=\langle\eps',\wh\delta\rangle.
\]
We fix the twist vectors
\[
 v_1=\eps,
 \qquad v_2=-\eps'.
\]
Thus $\gamma(v_1)=1$ and $\gamma(v_2)=-1$.
\item If $G=\langle T_1,T_2\rangle$, then
\[
 V^G=\Z\gamma,
 \qquad
 \sum_{g\in G}\im(g-I)=\langle\gamma,u,w\rangle,
\]
while
\[
 \Lambda^G=\Z\wh\delta,
 \qquad
 \sum_{g\in G}\im(A(g)-I)=\ker\gamma.
\]
In particular the coinvariants $V_G$ and $\Lambda_G$ are free of rank one.
The smallest $\langle A_1,A_2\rangle$-invariant subgroup of $\Lambda$
containing $\Lambda_{\mathrm{tor}}$ is $\ker\gamma$.
\item Set
\[
 q=uw+6\gamma\delta,
 \qquad
 \vol=\gamma uw\delta.
\]
Then
\[
 (\textstyle\bigwedge^kV)^G=
 \begin{cases}
 \Z,&k=0,\\
 \Z\gamma,&k=1,\\
 \Z q,&k=2,\\
 \Z\gamma uw,&k=3,\\
 \Z\vol,&k=4,
 \end{cases}
\]
and the coinvariants are
\[
 V_G=\Z[\delta],
 \quad
 (\textstyle\bigwedge^2V)_G=\Z[\gamma\delta],
 \quad [uw]=6[\gamma\delta],
 \quad
 (\textstyle\bigwedge^3V)_G=\Z[uw\delta],
 \quad
 (\textstyle\bigwedge^4V)_G=\Z[\vol].
\]
Finally
\[
\begin{aligned}
 V^{T_0}&=\langle\gamma,u\rangle,\\
 (\textstyle\bigwedge^2V)^{T_0}
 &=\langle\gamma u,\gamma\delta,uw,\gamma w-u\delta\rangle,\\
 (\textstyle\bigwedge^3V)^{T_0}
 &=\langle\gamma uw,\gamma u\delta\rangle.
\end{aligned}
\]
All displayed invariant subgroups are saturated, and all displayed
coinvariant groups are torsion-free.
\end{enumerate}
\end{proposition}

\begin{proof}
Direct multiplication gives $T_1^3=T_2^4=I$ and the displayed formula
for $T_0$.  The formulas for $A_1,A_2,M_0$ follow by taking inverse
transposes.  For
$\lambda=a\wh\gamma+b\wh u+c\wh w+d\wh\delta$,
\[
 (M_0-I)\lambda=-a\wh\delta+b\wh w,
\]
which proves (iii).  Solving $(A_j-I)\lambda=0$ gives
$b=2a,c=-4a$ for $j=1$, and $b=3a,c=-3a$ for $j=2$, proving (iv).

For (v), solving $T_jx=x$ gives
\[
 V^{T_1}=\langle\gamma,2u+w+3\delta\rangle,
 \qquad
 V^{T_2}=\langle\gamma,u+w+2\delta\rangle,
\]
whose intersection is $\Z\gamma$.  The images of $T_j-I$ contain
$u,w,2\gamma,3\gamma$, hence generate
$\langle\gamma,u,w\rangle$.  Dually, the images of $A_j-I$ generate
$\langle\wh u,\wh w,\wh\delta\rangle=\ker\gamma$.  The last assertion
follows because an invariant subgroup containing $\wh w$ contains
$A_1\wh w=\wh u-\wh w$, hence $\wh u$ and $\wh\delta$.

The exterior-power calculations in (vi) are given in
\cref{sec:matrix-tables}.  The Smith normal forms there have only unit
nonzero elementary divisors, which proves the asserted saturation.  The
$T_0$-invariants also follow directly from
$Nw=-u$, $N\delta=\gamma$, and $N\gamma=Nu=0$.
\end{proof}

\subsection{The orbifold curve}
Let $B=\PP^1$ with affine coordinate $t$, and put
\[
 p_1=0,
 \qquad p_2=1,
 \qquad p_0=\infty,
 \qquad t_c=t^{-1},
 \qquad B^\circ=B\setminus\{p_0,p_1,p_2\}.
\]
Let $\mathcal B$ be the orbifold with underlying curve $B$, cone orders
$3$ and $4$ at $p_1,p_2$, and a cusp at $p_0$.

\begin{proposition}[Uniformization]\label{prop:uniformization}
There are a copy $\HH_z$ of the upper half-plane, a Fuchsian group
$\Delta<\mathrm{PSL}_2(\R)$, and a holomorphic quotient map
\[
 \pi:\HH_z\longrightarrow B\setminus\{p_0\}
\]
with the following properties.
\begin{enumerate}[label=\textup{(\roman*)}]
\item $\HH_z/\Delta\cong B\setminus\{p_0\}$ as orbifolds.
\item There are $z_j\in\pi^{-1}(p_j)$ and generators $g_1,g_2$ such that
\[
 \Delta=\langle g_1,g_2\mid g_1^3=g_2^4=1\rangle,
 \qquad g_0=(g_1g_2)^{-1}
\]
is primitive parabolic, and
\[
 g_j'(z_j)=e^{-2\pi i/m_j},
 \qquad (m_1,m_2)=(3,4).
\]
\item For sufficiently small $r>0$, the component
$U_r\subset\pi^{-1}(0<|t_c|<r)$ accumulating at the fixed point of
$g_0$ is simply connected, invariant under $\langle g_0\rangle$, disjoint
from all its other $\Delta$-translates, and
\[
 U_r/\langle g_0\rangle\cong\{0<|t_c|<r\}.
\]
\end{enumerate}
\end{proposition}

\begin{proof}
This is the standard uniformization of the hyperbolic triangle orbifold of
signature $(0;3,4,\infty)$; see \cite[Chs.~7,10]{Beardon}.  Choose a hyperbolic triangle with angles
$\pi/3,\pi/4,0$ and take the orientation-preserving subgroup of its
reflection group.  If the vertices are ordered counterclockwise, the
inverses of the positive rotations about the finite vertices give
$g_1,g_2$ with the stated derivatives, while $(g_1g_2)^{-1}$ is the
primitive cusp generator.  The standard disjoint-horodisc theorem for a
cofinite Fuchsian group gives (iii).
\end{proof}

The homomorphism $\rho_V:\Delta\to\SL(V)$ is defined by
$\rho_V(g_j)=T_j$.  The monodromy on $\Lambda$ will be the contragredient
representation $g_j\mapsto A_j$, $g_0\mapsto M_0$.

